
Throughout this section, fix a category $\bcat$ which admits all finite limits
and a symmetric monoidal functor
\[ \D{-} \colon \Span{\bcat} \too \Prlst,\] where $\Span{\cC}$ denotes the 1-category
of spans in $\bcat$ equipped with the cartesian monoidal structure.
We write $\D{\pt}=\cD$ and note that $\D{\pt}$ refines to an object
$\D{\pt} \in \CAlg(\Prlst)$ and $\D{-}$ naturally refines to a symmetric monoidal functor
\[ \D{-} \colon \Span{\bcat} \too \Mod_{\cD}(\Prlst)\] for which we abusively use
the same notation. We also make the standing assumption that $\cD$ is compactly
generated and \tdef{rigid}, that is, an object of $\cD$ is compact if and only
if it is dualizable.
For a map $f\colon X\to Y$ in $\bcat$ we denote by
\[ f^{\ast}\colon \D{Y} \too \D{X}\] the $\D{\pt}$-linear left adjoint obtained by
applying $\D{-}$ to the span
\[\begin{tikzcd}
    & X \\
    Y && X \arrow["f"', from=1-2, to=2-1] \arrow["\id", from=1-2, to=2-3].
  \end{tikzcd}\] Conversely, we denote by
\[ f_{!}\colon \D{X} \too \D{Y}\] the $\D{\pt}$-linear left adjoint obtained by
applying $\D{-}$ to the span
\[\begin{tikzcd}
    & X \\
    X && Y \arrow["\id"', from=1-2, to=2-1] \arrow["f", from=1-2, to=2-3],
  \end{tikzcd}\] and denote the induced adjoints as
\[ f^{\ast}\dashv f_{\ast} \qquad f_{!} \dashv f^{!}.\]
\begin{defn}
  For any $X\in \bcat$ and $\cE\in \D{X}$, we write
  $\Gamma^{c}(X;\cE):=X_{!}\cE \in \D{\pt}$ and $\Gamma(X;\cE):= X_{\ast}\cE$. Moreover, for
  any map $f\colon X\to Y$ in $\bcat$ we write $\omega_{X/Y}:= f^{!}\1_{\D{Y}}$ and call
  this the \tdef{relative dualizing sheaf} of $f$. For $f=X\to \pt$, being the
  terminal map, we also write $\omega_{X}:=X^{!}\1$ and call this the (absolute)
  \tdef{dualizing sheaf} of $X$.
\end{defn}

\begin{rmk}
  The category $\D{X}$ inherits a natural presentably symmetric monoidal
  structure via the functors
  \[ X^{\ast}\colon \D{\pt} \too \D{X}\]
  \[ \Delta^{\ast}\colon \D{X \times X} \simeq \D{X} \otimes \D{X} \too \D{X}\]
\end{rmk}
Since $\D{X}$ is presentably monoidal, the functor
\[ \cE \otimes - \colon \D{X} \too \D{X},\] admits a right adjoint, which we denote by
\[ \Hom_{\D{X}}(\cE, -) \colon \D{X} \too \D{X}.\] Such that, for any $\cF\in \D{X}$ we
have an equivalence
\[ \Gamma(X;\Hom(\cE,\cF))=X_{\ast} \Hom_{\D{X}}(\cE, \cF) \simeq \map_{\D{X}}(\cE,\cF) \in \D{\pt}. \]

\begin{prop}\label{prop: projection formula}
  Let $f\colon X\to Y$ be a map in $\bcat$. Then for any $\cE, \cE^{\prime}\in \D{X}$ and
  $\cF\in \D{Y}$, we have the following natural equivalences
  \[ f_{!}(\cE \otimes f^{\ast}\cF) \simeq f_{!}(\cE)\otimes \cF\]
  \[\Hom_{\D{Y}}(f_{!}\cE, \cF)\simeq f_{\ast}\Hom_{\D{X}}(\cE, f^{!}\cF) \]
  \[ f^{!}\Hom_{\D{X}}(\cE, \cE^{\prime})\simeq \Hom_{\D{Y}}(f^{\ast}\cE, f^{!}\cE^{\prime} ).\]
\end{prop}
\begin{cons}
  Since any $X\in \Span{\bcat}$ is canonically self-dual, we also obtain a natural
  equivalence
  \[ \VD \colon \D{X} \rar{\sim} \D{X}^{\vee}, \quad \cE \mapsto X_{!}(- \otimes \cE)\] associating to
  each $\cE\in \D{X}$ the continuous functor given by the composite
  \[ \D{X} \rar{\pi_{1}^{\ast}} \D{X} \otimes \D{X} \rar{ \id \otimes \cE } \D{X} \otimes \D{X} \rar{\Delta^{\ast}} \D{X} \rar{X_{!}} \D{X}.\]
  This comes together with, for any map $f\colon X\to Y$, identifications
  \[ f^{\ast} \simeq \VD\circ~(f_{!})^{\vee}\circ \VD \]
  \[ f_{!} \simeq \VD \circ~(f^{\ast})^{\vee}\circ \VD.\] In particular, composing with the
  internal dual in $\D{X}^{\vee}$, we obtain a functor
  \[ \Dl\colon \D{X} \rar{\VD} \D{X}^{\vee} \rar{(\Hom_{\D{X}}(-,\1))^{\vee}} (\D{X}^{\vee})^{\op} \rar{ (\VD^{-1})^{\op}} \D{X}^{\op}.\]
  Unwinding the definitions, we see that $\Dl$ is given by
  \[ \Dl = \Hom_{\D{X}}(-, \omega_{X}) = \Hom_{\D{X}}(-, X^{!}\1). \]
\end{cons}

\begin{defn}
  For any $X\in \bcat$, we say that $\cE\in \D{X}$ is \tdef{reflexive}, if the
  canonical map $\Dl^{2}(\cE) \to \cE$ is an equivalence.
\end{defn}

\begin{lem}\label{lem:push pull dual}
  Let $f\colon X\to Y$ be a map in $\bcat$ and $\cE\in \D{X}$ and $\cF\in \D{Y}$. Then, we
  have natural equivalences
  \[ \Dl(f_{!}\cE)\simeq f_{\ast}\Dl(\cE) \in \D{Y}\]
  \[ f^{!}\Dl(\cF) \simeq \Dl(f^{\ast}\cF) \in \D{X}.\]
\end{lem}
\begin{proof}
  Indeed, the first equivalence is given by the composite
  \[ \Dl (f_{!}\cE) = \Hom(f_{!}\cE, \omega_{Y}) \simeq f_{\ast}\Hom(\cE, f^{!}\omega_{Y}) \simeq f_{\ast}\Hom(\cE, \omega_{X}) \simeq f_{\ast}\Dl(\cE), \]
  while the second is obtained as
  \[ f^{!}\Dl(\cF)= f^{!}\Hom(\cF, \omega_{Y}) \simeq \Hom(f^{\ast}\cE, \omega_{X}) =\Dl(f^{\ast}\cF),\]
  proving our claim.
\end{proof}

Combining the two equivalences, we learn that for a sheaf $\cF\in \D{Y}$ we have
equivalences
\[ \Gamma(X; f^{!}(\Dl(\cF)))\simeq \Gamma(X;\Dl(f^{\ast}\cF))= X_{\ast} \Dl(f^{\ast}\cF) \simeq \Dl(X_{!}f^{\ast}\cF) \simeq \Gamma^{c}(X;f^{\ast}\cF)^{\vee}.\]
In particular, for $Y=X$ and $f=\id$ we have:
\[ (X_{!}\cE)^{\vee}=\Gamma^{c}(X;\cE)^{\vee} \simeq \Gamma(X;\Dl(\cE))=X_{\ast}\Dl(\cE).\] Put
differently, this means that the functor $\D{X} \to \D{\pt}$ which takes $\cE$ to
$\Gamma^{c}(X;\cE)^{\vee}$ is represented by the dualizing sheaf $\omega_{X}$. If $X$ is
additionally assumed to be proper, i.e.~$X_{!}=X_{\ast}$, we may iterate this
reasoning to see that
\[ \Gamma(X;\Dl^{2}(\cE)) \simeq (\Gamma(X;\cE)^{\vee})^{\vee}.\]
\begin{cor}\label{cor:coherent push}
  For a proper map $f\colon X\to Y\in \bcat$, the functor $f_{!}$ preserves reflexive
  sheaves.
\end{cor}
\begin{proof}
  Since $f$ is proper, we have $f_{!}=f_{\pt}$ and so applying \cref{lem:push pull dual} twice
  we get for any reflexive $\cE\in \D{X}$ equivalences
  \[ \Dl^{2}(f_{!}\cE) \simeq \Dl(f_{!}\Dl(\cE)) \simeq f_{!}\Dl^{2}(\cE) \simeq f_{!}\cE.\]
  Thus, $f_{!}\cE\in \D{Y}$ is by definition reflexive.
\end{proof}

Notice that, for any rfelexive $\cF,\cE\in \D{X}$ we have
\[\Hom(\cF, \Dl(\cE)) \simeq \Hom(\cF \otimes \cE, \omega_{X})
  \simeq \Dl(\cF \otimes \cE).\] Replacing $\cF$ with $\Dl(\cF)$, we learn that we can
compute the internal mapping sheaf between reflexive sheaves sheaves as the tensor
product
\[ \Hom(\cF, \cE) \simeq \Hom(\cF, \Dl(\Dl(\cE))) \simeq \Dl(\cF \otimes \Dl(\cE)) \in \D{X}.\]


\begin{lem}\label{lem:suave conditions}
  Let $f\colon X\to Y$ be a proper map in $\bcat$ and let $\cE\in \D{X}$. The
  following statements are equivalent:
  \begin{enumerate}
    \item The functor $f^{\ast}(-)\otimes \cE$ has a left adjoint
    \item The functor $f_{!}(\cE \otimes -)$ has a continuous right adjoint
    \item The functor $\Hom(\cE, f^{!}(-))$ has a right adjoint
  \end{enumerate}
  Moreover, in this case we have $\Dl^{2}(\cE)\simeq \cE$ and the adjunctions are
  given by
  \[ f_{!}(\Dl(\cE)\otimes - ) \dashv f^{\ast}(-)\otimes\cE \simeq \Hom(\Dl(\cE), f^{!}(-))\] and
  \[ \Hom(\cE, f^{!}(-)) \simeq f^{\ast}(-)\otimes \Dl(\cE) \dashv f_{!}(\Dl(\cE)^{\vee} \otimes - )\]
\end{lem}
\begin{proof}
  Note that $f_{!}(\cE\otimes -)$ always admits a right adjoint given by
  $\Hom(\cE, f^{!}(-))$, so conditions (2) and (3) are clearly equivalent.
\end{proof}

\begin{defn}
  Let $f\colon X\to Y$ be a map in $\bcat$. If $\cE\in \D{X}$ satisfies any of the
  equivalent conditions of \cref{lem:suave conditions}, we call $\cE$ \tdef{$f$-suave} or just
  \tdef{suave} if $Y=\pt$ and $f$ is the terminal map. We write $\Ds{X}$ for the
  full subcategory of $\D{X}$ spanned by the suave sheaves on $X$. Moreover, we
  call $X\in \bcat$ suave if the unit $\1_{X}\in \D{X}$ is suave and \tdef{smooth}
  if additionally $\omega_{X}=X^{!}\1$ is dualizable.
\end{defn}

In particular, suaveness ensures that the cohomology $\Gamma(X)$ is dualizable in $\D{\pt}$.


\begin{prop}\label{prop: sections dualizable}
  Let $X\in \bcat$ be proper and $\cE\in \D{X}$ be suave. If $\1_{\cD}\in \cD^{\omega}$,
  the global sections $\Gamma(X;\cE)$ are a dualizable object of $\D{\pt}$.
\end{prop}
\begin{proof}
  Since $X$ is proper, the functor $X^{\ast}$ is strongly continuous and hence
  preserves compact objects, hence the unit $X^{\ast}\1=\1_{X}\in \D{X}$ is
  compact. Moreover, since $\cE$ is suave, the functor $X_{!}(-\otimes \cE)$
  is strongly continuous and thus $\Gamma(X;\cE)=X_{!}(\1_{X}\otimes \cE)\in \D{X}$
  is compact and hence dualizable, since $\D{\pt}$ is locally rigid..
\end{proof}

For $X\in \bcat$, denote by $\pi_{i}X\times X\to X$ the canonical projections and by
$\Delta\colon X \to X\times X$ the diagonal. The \tdef{external tensor product}
of two sheaves $\cE,\cF\in \D{X}$ is defined to be
\[ \cE\boxtimes \cF := \pi_{1}^{\ast}\cE \otimes \pi_{2}^{\ast}\cF \in \D{X\times X}.\]
Note that, we have a natural equivalence
\[ \Delta^{\ast}(\cE\boxtimes \cF)\simeq \cE\otimes \cF\in \D{X}.\]

\begin{prop}\label{prop: external hom}
Let $\cE,\cF\in \D{X}$ be suave and denote by $\pi_{i}\colon X\times X \to X$ the two projections.
Then, we have a natural equivalence
\[ \pi_{1}^{\ast}\Dl(\cE)\otimes \pi_{2}^{\ast}\cF \rar{\sim} \Hom(\pi_{1}^{\ast}\cE, \pi_{2}^{!}\cF)\in \D{X\times X}\]
\end{prop}

\begin{cor}
  If $\cE,\cF\in \D{X}$ are suave, we have a natural equivalence
  \[ \Delta^{!}(\Dl(\cE)\boxtimes \cF)\simeq\Hom(\cE, \cF) \in \D{X} .\]
\end{cor}
\begin{proof}
  Follows by applying $\Delta^{!}$ to the equivalence of \cref{prop: external hom} and
  observing that we have
  \[ \Delta^{!}\Hom(\pi_{1}^{\ast}\cE, \pi_{2}^{!}\cF)
    \simeq \Hom(\Delta^{\ast}\pi_{!}^{\ast}\cE, \Delta^{!}\pi_{2}^{!}\cF)
  \simeq \Hom(\cE, \cF)\]
by \cref{prop: projection formula} and the fact that $\pi_{i}\circ \Delta=\id$.
\end{proof}
% \end{prop}
% \begin{proof}
%   The ``if'' direction is clear by Lemma~\ref{lem:suave conditions}. For the only if direction first
%   note that $\cE$ is suave if and only if $\Dl(\cE)$ is suave. Moreover, if
%   the functor $X^{\ast}(-)\otimes \Dl(\cE)$ admits a left adjoint, it is necessarily
%   given by $X_{!}(\Dl^{2}(\cE) \otimes -)$. However, the latter functor always
%   admits a right adjoint, which is given by $\Hom(\cE,X^{!}(-))$. Since $X$ is
%   suave we have $X^{!}(-)\simeq X^{\ast}(-)\otimes \omega_{X}$ and hence $\Dl(\cE)$ is suave if
%   and only if we have an equivalence
%   \[ X^{\ast}(-)\otimes \Dl(\cE) \simeq \Hom(\cE,X^{\ast}(-)\otimes \omega_{X}) \] Since $X$ was assumed
%   to be smooth i.e.~$\omega_{X}$ is dualizable, we obtain an equivalences
%   \[ \omega_{X}\otimes \cE^{\vee}\simeq \omega_{X} \otimes \Hom(\cE, \1_{X}) \simeq \Hom(\cE, \omega_{X}) = \Dl(\cE)\]
%   as claimed.
% \end{proof}


% \begin{obs}\label{obs: diagonal comparison}
%   Let $X\in \bcat$ be separated, then there is a canonical comparison map
%   $\delta\colon \Delta^{!}\to \Delta^{\ast}$ given by the composite
%   \[ \delta\colon \Delta^{!}\simeq \Delta^{\ast}\Delta_{\ast}\Delta^{!} \simeq \Delta^{\ast}\Delta_{!}\Delta^{!}\too \Delta^{\ast},\] where the last map
%   is the unit of the adjunction $\Delta_{!}\dashv \Delta^{!}$.
% \end{obs}

% \begin{prop}
%   Let $X\in \bcat$ be separated and suave, denote the diagonal by $\Delta\colon X\to X\times X$ and
%   the two projections by $\pi_{i}\colon X\times X \to X$. For any $\cE,\cF\in \Ds{X}$, the
%   external tensor product
%   \[ \cE \boxtimes \cF:= \pi_{1}^{\ast}\cE \otimes \pi_{2}^{\ast} \cF \in \D{X\times X}\] is suave with Verdier
%   dual given by
%   \[ \Dl(\cE \boxtimes \cF)\simeq \Dl(\cE)\boxtimes \Dl(\cF)\in \D{X\times X}.\] In particular, we can
%   compute the mapping space as the exceptional inverse image
%   \[ \Hom(\cF, \cE) \simeq \Delta^{!}(\cE \boxtimes \Dl(\cF)) \in \D{X}.\]
% \end{prop}

% \begin{defn}\label{def: character}
%   Let $X$ be separated and suave. For any $\cE\in \Ds{X}$ we define the
%   \tdef{character} of $\cE$ as the composite map
%   \[ \chi^{\cE}\colon\1_{X}\rar{\id} \Hom(\cE, \cE)\simeq \Delta^{!}(\cE \boxtimes \Dl(\cE)) \rar{\delta} \Delta^{\ast}(\cE \boxtimes \Dl(\cE))\simeq \cE \otimes \Dl(\cE) \rar{\ev} \omega_{X}.\]
%   Equivalently, under the adjunction $X^{\ast}\dashv X_{\ast}$, the map $\chi^{\cE}$ defines
%   an element in the global sections $\chi^{\cE}\colon \1 \to \Gamma(X;\omega_{X})$, i.e.~a class in
%   the homology of $X$.
% \end{defn}
% \begin{prop}
%   For $X\in \bcat$ smooth with proper diagonal, the category $\Dc{X}$ inherits
%   the structure of a small, idempotent complete, stable, symmetric monoidal
%   category from $\D{X}$.
% \end{prop}
% \begin{proof}
%   Since $\Dl=\Hom_{\D{X}}(-,\omega_{X})$, the functor $\Dl^{2}$ commutes with
%   finite limits and colimits. Moreover, since any functor preserves retracts,
%   the first claim follows. Now, for any $\cE \in \D{X}^{\dual}$ we have
%   equivalences
%   \[ \Dl^{2}\cE \simeq \Dl(\Hom_{\D{X}}(\cE, \omega_{X})) \simeq \Hom_{\D{X}}(\omega_{X}\otimes \cE^{\vee},\omega_{X}) \simeq \omega_{X} \otimes (\omega_{X} \otimes \cE^{\vee})^{\vee} \simeq \cE,\]
%   and thus $\D{X}^{\dual}\subseteq \Dc{X}$, as claimed. Smallness should follows since
%   $\Dl$ induces an equivalence $\Dc{X} \simeq \Dc{X}^{\op}$ and the final claim is
%   clear since $\D{X}$ is stable and idempotent complete.
% \end{proof}
% \Fnote{The next bit is wrong cause of monoidality}

% \begin{cor}
%   Let $X\in \bcat$ be smooth and proper and let $\cE,\cF\in \D{X}$ be coherent.
%   Then $X_{!}\cE= \Gamma^{c}(X;\cE)$ and $\map(\cE ,\cF)$ are dualizable objects in
%   $\D{\pt}$.
% \end{cor}
% \begin{proof}
%   Indeed, we have that
%   \[ \map(\cE,\cF) = X_{!}\Hom(\cE,\cF)\simeq X_{!}\Dl((\cF\otimes \Dl(\cE)) = \Gamma^{c}(X;\Dl(\cF \otimes \Dl(\cE))),\]
%   and hence both claims follow from the previous lemma, using that
%   $\Dc{\pt}=\D{\pt}^{\dual}$ holds by definition.
% \end{proof}

% Let $\Cat^{\perf}_{\D{\pt}^{\dual}}$ denote the category of small, stable,
% idempotent complete, $\D{\pt}^{\dual}$-linear categories with maps given by
% exact, $\D{\pt}$-linear functors.

% \begin{prop}
%   For any smooth and proper $X\in \bcat$, the functor
%   \[ \Dc{X}^{\op} \too \Fun^{\rm{ex}}_{\Dc{\pt}}(\Dc{X},\Dc{\pt}), \quad \cE \mapsto X_{!}(\Hom(\cE, -)) \simeq X_{!}( - \otimes \Dl(\cE))\]
%   is an equivalence. Moreover, the category $\Dc{X}$ is dualizable in
%   $\Cat^{\perf}_{\D{\pt}^{\dual}}$ with dual $\Dc{X}^{\vee} = \Dc{X}^{\op}$ and
%   evaluation given by the enriched Yoneda functor
%   \[ X_{!}\Hom(-,-)\colon \Dc{X} \times \Dc{X}^{\op} \too \D{\pt}^{\dual}\] and
%   coevaluation given by the functor
%   \[ (X^{\ast}, \Dl(X^{\ast}))\colon \D{\pt}^{\dual} \too \Dc{X}\times \Dc{X}^{\op}.\]
% \end{prop}

% This allows us to fix the defect that the pullback functor $f^{\ast}$ has no
% reason to preserve coherent sheaves.
% \begin{cons}
%   Let $f\colon X\to Y \in \bcat$ be a proper map of proper, smooth spaces. We define
%   the \tdef{coherent pullback} along $f$ as the functor obtained by the
%   composition
%   \[ f^{\flat}\colon\Dc{Y} \rar{\sim} \Dc{Y}^{\vee} \rar{f_{!}^{\vee}} \Dc{X}^{\vee} \rar{\sim} \Dc{X}.\]
% \end{cons}

% \begin{defn}
%   Let $X\in \bcat$ be smooth and proper. We write $\IcD{X}$ for the Ind-category
%   of $\Dc{X}$ and call this the category of \tdef{ind-coherent} sheaves on
%   $X$.
% \end{defn}

% Write $\bcat^{sp}$ for the subcategory of $\bcat$ spanned by the smooth and
% proper objects with proper maps between them.

% \begin{prop}
%   For $X\in \bcat^{sp}$, the category $\IcD{X}$ is stable, presentable,
%   compactly generated and inherits the structure of a presentably symmetric
%   monoidal $\D{\pt}$-linear category. Moreover, $\IcD{X}$ is fully dualizable
%   and canonically self-dual.
% \end{prop}

% For $f\colon X\to Y \in \bcat^{sp}$, the functors $f^{\flat}$ and $f_{\ast}$ induce continuous
% functors on Ind-categories
% \[ f_{!}\colon \IcD{X} \too \IcD{Y} \qquad f^{\flat} \colon \IcD{Y} \too \IcD{X}.\]
% \begin{thm}
%   The functors above refine to a symmetric monoidal functor
%   \[ \IcD{-}\colon \Span{\bcat^{sp}} \too \Prlst\] i.e.~a six functor formalism.
% \end{thm}
% \subsection{Ramification}

% Let $X\in \bcat$ be a space. We want to measure ``singularities'' of $X$.

% \begin{defn}
% For a sheaf $\cE\in \D{X}$ we define its \tdef{vanishing set} to be the set of points $V(\cE)\subseteq X$
% at which the stalk $x^{\ast} \cE$ vanishes. The \tdef{support} of $\cE$ is define to be
% $S(\cE):= X\setminus V(\cE)$.
% \end{defn}

% \begin{defn}
%   We define a sheaf $\cR_{X} \in \D{X}$ called the \tdef{ramification sheaf} of $X$
%   as the fiber of the evaluation map
%   \[ \cR_{X}:= \fib \left( \omega_{X} \otimes \omega_{X}^{\vee} \too \1_{X} \right)\]
%   We say that a point $x\in X$ is \tdef{singular} if $x\in S(\cR_{X})$ and \tdef{regular}
%   else. We call $X$ regular if all of its points are regular.
% \end{defn}

% Note that $X$ is regular if and only if $\omega_{X}$ is \textit{invertible}.

% \begin{ex}
%   Let $X$ be a graph and let $x\in X$ be a point. Then the support of $\cR_{X}$ is given by those
%   vertices which have more than 2 incoming edges.
% \end{ex}

% \begin{ex}
%   Let $M$ be a manifold of dimension $n$. Then $\omega_{X}$ is locally isomorphic to the constant
%   sheaf $\S[n]$ and hence invertible, so $M$ is regular. In fact, being regular is equivalent
%   to being a homology manifold.
% \end{ex}

%%% Local Variables:
%%% TeX-master: "main"
%%% End:
